\BeleskeID{06-s005}
% element: 06-s005:e01
% element: 06-s005:d01
Da bi se binarna vrednost prikazala decimalnim ciframa u tekstu ili na ekranu, može se ponavljati celobrojno deljenje sa 10. Ostatak je cifra najmanje težine, a količnik se ponovo deli sa 10 dok se ne izdvoje sve cifre.

Za račun $59+13$ sabiraju se kodovi jedinica $1001+0011=1100$, odnosno binarna vrednost 12. Ona nije dopuštena decimalna cifra u kodu 8421. Dodavanjem $0110$ dobija se $1\,0010$: donja četiri bita predstavljaju cifru 2, a jedinica se prenosi u kolonu desetica. Desetice zato postaju $5+1+1=7$.

Korekcija šest razlikuje binarnu osnovu 16 od decimalne osnove 10. Vrednost koja prelazi decimalnu cifru mora dati prenos deset, dok četvorobitno sabiranje prirodno daje prenos šesnaest. Zato je potrebna razlika $16-10=6$.
